%=====================================================================
\section{Psychrometric foundations}\label{sec:psy:top}
\label{sec:psy:foundations}

Everything the plant does to air --- heating it in the regeneration coil, drying it in the desiccant
wheel, cooling it before it enters the hatchery --- is bookkeeping on two numbers per stream. This
section derives those numbers and the relations that move them around, exactly as implemented in the
notebook (\texttt{dbhe\_master.ipynb}, v10.6, psychrometric block of the desiccant-wheel cell). The
equations are elementary and unforgiving: the whole latent story of this plant lives in the third
decimal place of a humidity ratio.

\subsection{Why the humidity ratio, and why ``per kilogram of dry air''}
\label{sec:psy:why-w}

Moist air is a mixture of dry air and water vapour. As a stream travels around either loop, water is
added (desorption in the wheel, egg metabolism, infiltration) and removed (adsorption in the wheel),
so the \emph{total} moist-air mass flow is not conserved across a component. The dry-air mass flow
is: nothing here creates or destroys dry air. That single fact fixes the convention. Define the
humidity ratio
\begin{equation}
w \;\equiv\; \frac{m_{\mathrm{v}}}{m_{\mathrm{a}}}
\;=\; \frac{\text{mass of water vapour}}{\text{mass of dry air}}
\qquad [\si{\kilogram\per\kilogram}] , \label{eq:psy:wdef}
\end{equation}
and carry every extensive property \emph{per kilogram of dry air}. A water balance on any component
is then literally ``$\md w$ in equals $\md w$ out'', with the same constant $\md$ on both sides: the
bookkeeping is linear and cannot drift. Had we used relative humidity, or moist-air mass, every
balance would carry a denominator that changes as the process runs, and small errors would compound
around the loop. This is why the notebook passes $(T,w)$ pairs --- never $(T,\RH)$ --- between
components; relative humidity is computed only where a human or a warning message has to read it.

\subsection{Humidity ratio from partial pressures}
\label{sec:psy:partial}

Treat dry air and vapour as ideal gases sharing a volume $V$ at temperature $T$, with partial
pressures summing to the total by Dalton's law, $\patm = p_{\mathrm{a}} + p_{\mathrm{v}}$. Then,
using $R_i = R_u/M_i$,
\begin{equation}
w \;=\; \frac{m_{\mathrm{v}}}{m_{\mathrm{a}}}
\;=\; \frac{p_{\mathrm{v}} V/(R_{\mathrm{v}} T)}{p_{\mathrm{a}} V/(R_{\mathrm{a}} T)}
\;=\; \frac{p_{\mathrm{v}}}{p_{\mathrm{a}}}\,\frac{M_{\mathrm{v}}}{M_{\mathrm{a}}} .
\label{eq:psy:w_ideal}
\end{equation}
With $M_{\mathrm{v}} = \SI{18.015}{\gram\per\mole}$ and $M_{\mathrm{a}} =
\SI{28.966}{\gram\per\mole}$ the molar-mass ratio is $0.6219$, which is the constant hard-coded (as
$0.622$) in every psychrometric routine of the notebook:
\begin{equation}
\boxed{\;w \;=\; 0.622\,\frac{p_{\mathrm{v}}}{\patm - p_{\mathrm{v}}}\;} \label{eq:psy:w_from_pv}
\end{equation}

Two things to notice. First, $w$ depends on the \emph{total} pressure --- which is why $\patm =
\SI{101325}{\pascal}$ is a declared configuration constant and every psychrometric call is routed
through that one symbol: at another altitude the same $T$ and $\RH$ give a different $w$. Second,
the relation is non-linear in $p_{\mathrm{v}}$, because the denominator shrinks as the air gets
wetter; at the humidities this plant works in (\SIrange{20}{22}{\gram\per\kilogram}) that curvature
is a few percent, not negligible.

\subsection{Saturation pressure and relative humidity}
\label{sec:psy:psat}

Relative humidity is the ratio of the actual vapour partial pressure to the saturation pressure at
the \emph{same} temperature,
\begin{equation}
\RH \;\equiv\; \frac{p_{\mathrm{v}}}{\psat(T)} , \qquad\text{equivalently}\qquad
p_{\mathrm{v}} \;=\; \RH \cdot \psat(T) . \label{eq:psy:RHdef}
\end{equation}

The code does not call a steam table for $\psat$. It uses a single-line Magnus-type correlation,
with $T$ in \si{\celsius} and $\psat$ in \si{\pascal}:
\begin{equation}
\boxed{\;\psat(T) \;=\; 610.94\, \exp\!\left(\frac{17.625\,T}{T + 243.04}\right)\;}
\label{eq:psy:psat}
\end{equation}
These are the Alduchov--Eskridge coefficients for saturation over liquid water, accurate to better
than \SI{0.4}{\percent} over \SIrange{-40}{50}{\celsius}. That covers the process loop comfortably,
but the regeneration air runs at \SIrange{60}{100}{\celsius}, outside the fitted range, and
\eqref{eq:psy:psat} is still evaluated there by the wheel's relative-humidity diagnostic. The cost
is confined to that diagnostic: $\psat$ never enters a mass or energy balance, it only converts $w$
to $\RH$ for display and regime checking.

Substituting \eqref{eq:psy:RHdef} and \eqref{eq:psy:psat} into \eqref{eq:psy:w_from_pv} gives the
routine \texttt{w\_from\_RH}$(T,\RH)$ used throughout:
\begin{equation}
w(T,\RH) \;=\; 0.622\,\frac{\RH\,\psat(T)}{\patm - \RH\,\psat(T)} . \label{eq:psy:w_from_RH}
\end{equation}
The implementation clips $\RH$ to $[0,1]$ first, so a caller cannot request a supersaturated state
through this door --- the state simply saturates. That is a guard, not physics: it hides an
inconsistent request rather than reporting it.

Inverting \eqref{eq:psy:w_from_pv} for $p_{\mathrm{v}}$ and dividing by $\psat$ gives the reverse
map:
\begin{equation}
\boxed{\;\RH(T,w) \;=\; \frac{1}{\psat(T)}\,\frac{w\,\patm}{0.622 + w}\;} \label{eq:psy:RH_from_w}
\end{equation}
There is no single named \texttt{RH\_from\_w} helper in v10.6; this same expression is written out
three times (in the wheel regime diagnostic, in the results dashboard, and inline in the load
interface), algebraically identical in all three.

\paragraph{The consequence that drives this project.}
$\psat$ grows nearly exponentially with $T$, so a fixed relative humidity at a high temperature is a
very large absolute moisture content. The hatchery setpoint of \SI{37}{\celsius} at
\SI{50}{\percent}~RH is not a dry condition: from \eqref{eq:psy:psat}, $\psat(37) =
\SI{6271}{\pascal}$, so $p_{\mathrm{v}} = \SI{3136}{\pascal}$ and $w =
\SI{19.86}{\gram\per\kilogram}$. Whether the plant must dehumidify or humidify is settled entirely
by the sign of $w_{\mathrm{oa}} - w_{\mathrm{pc}}$ --- a comparison of humidity ratios, never of
relative humidities.

\subsection{Moist-air enthalpy as coded}
\label{sec:psy:enthalpy}

Take dry air at \SI{0}{\celsius} and saturated liquid water at \SI{0}{\celsius} as reference states.
The enthalpy per kilogram of dry air is the dry-air sensible term plus the vapour term (latent heat
at the reference, plus the vapour's own sensible heat):
\begin{equation}
h(T,w) \;=\; \underbrace{c_{p,\mathrm{a}}\,T}_{\text{dry air}}
\;+\; w\,\underbrace{\bigl(h_{\mathrm{fg},0} + c_{p,\mathrm{v}}\,T\bigr)}_{\text{vapour}} ,
\label{eq:psy:h_general}
\end{equation}
and the notebook fixes the three constants at $c_{p,\mathrm{a}} =
\SI{1.006}{\kilo\joule\per\kilogram\per\kelvin}$, $h_{\mathrm{fg},0} =
\SI{2501}{\kilo\joule\per\kilogram}$, $c_{p,\mathrm{v}} =
\SI{1.86}{\kilo\joule\per\kilogram\per\kelvin}$:
\begin{equation}
\boxed{\;h(T,w) \;=\; 1.006\,T + w\,\bigl(2501 + 1.86\,T\bigr)
\qquad [\si{\kilo\joule\per\kilogram}\ \text{dry air}]\;} \label{eq:psy:h_coded}
\end{equation}
Both specific heats are constants --- no temperature or pressure dependence. Over
\SIrange{0}{100}{\celsius} the true $c_{p,\mathrm{a}}$ moves by well under \SI{1}{\percent}, so the
cost is small, and as the next subsection shows it is what makes mixing exactly linear. Because
mixing delivers $(h,w)$ rather than $(T,w)$ we need the inverse, and solving \eqref{eq:psy:h_coded}
for $T$ is a one-line algebraic step precisely because the constants are constant:
\begin{equation}
h = 1.006\,T + 2501\,w + 1.86\,w\,T \quad\Longrightarrow\quad
\boxed{\;T(h,w) \;=\; \frac{h - 2501\,w}{1.006 + 1.86\,w}\;} \label{eq:psy:T_from_hw}
\end{equation}
This is \texttt{T\_from\_hw}. Note what the denominator is: the moist-air specific heat per kilogram
of dry air,
\begin{equation}
c_{p,\mathrm{m}}(w) \;=\; 1.006 + 1.86\,w \qquad [\si{\kilo\joule\per\kilogram\per\kelvin}] .
\label{eq:psy:cp_moist}
\end{equation}
Equation~\eqref{eq:psy:cp_moist} reappears in that exact form in the air-to-air cooler (at the
stream's own $w$), in the building load calculation (at the mean $\tfrac12(w_{\mathrm{pd}} +
w_{\mathrm{pc}})$), and in the generic load interface (at the room humidity). The regeneration coil
is the exception: it uses the constant dry-air value \SI{1006}{\joule\per\kilogram\per\kelvin} on
the air side and drops the $1.86\,w$ term entirely. At the regeneration inlet humidity, about
\SI{20}{\gram\per\kilogram}, that understates the air-side capacity rate --- and hence the coil duty
--- by roughly \SI{3.7}{\percent}. It is a bounded inconsistency, not a justified simplification.

A second note: \eqref{eq:psy:h_coded} and the building latent load use $h_{\mathrm{fg},0} =
\SI{2501}{\kilo\joule\per\kilogram}$, while the load generators and the load interface convert
moisture rate to latent power with a separate configuration constant $\hfg =
\SI{2.5}{\mega\joule\per\kilogram}$. They differ by \SI{0.04}{\percent} and no result depends on it,
but a reader comparing $\ql$ against an enthalpy difference will find a residual that is arithmetic,
not physics.

\subsection{The adiabatic mixing rule}
\label{sec:psy:mixing}

Every junction in both air loops is this one derivation, so it is worth doing carefully once. Two
streams with dry-air mass flows $\md_1$ and $\md_2$, at states $(T_1,w_1)$ and $(T_2,w_2)$, mix
adiabatically at constant pressure into one stream $\md_3$ at $(T_3,w_3)$. Three balances close it.
\emph{Dry air}, which nothing creates or destroys:
\begin{equation}
\md_3 \;=\; \md_1 + \md_2 . \label{eq:psy:mix_air}
\end{equation}
\emph{Water}, since stream $i$ carries $\md_i w_i$ of vapour:
\begin{equation}
\md_1 w_1 + \md_2 w_2 = \md_3 w_3 \quad\Longrightarrow\quad
w_3 = \frac{\md_1 w_1 + \md_2 w_2}{\md_1 + \md_2} . \label{eq:psy:mix_w}
\end{equation}
\emph{Energy}: with no work and no heat crossing the boundary, enthalpy is conserved:
\begin{equation}
\md_1 h_1 + \md_2 h_2 = \md_3 h_3 \quad\Longrightarrow\quad
h_3 = \frac{\md_1 h_1 + \md_2 h_2}{\md_1 + \md_2} . \label{eq:psy:mix_h}
\end{equation}
The mixed temperature then follows from \eqref{eq:psy:T_from_hw}, $T_3 = T(h_3,w_3)$. Mixing is done
in $(w,h)$ and converted back to $T$ at the end, never the other way round.

\paragraph{Why the temperature is not mass-averaged.}
It is tempting to write $T_3 = (\md_1 T_1 + \md_2 T_2)/\md_3$. In general this is wrong, and the
reason is instructive. Substituting \eqref{eq:psy:h_general} into \eqref{eq:psy:mix_h},
\begin{equation*}
c_{p,\mathrm{a}} T_3 + h_{\mathrm{fg},0} w_3 + c_{p,\mathrm{v}} w_3 T_3
\;=\; \frac{\md_1\bigl(c_{p,\mathrm{a}} T_1 + h_{\mathrm{fg},0} w_1
+ c_{p,\mathrm{v}} w_1 T_1\bigr) + \md_2(\cdots)}{\md_3} ,
\end{equation*}
the terms linear in $T$ and in $w$ separate cleanly and by themselves would give a mass-averaged
$T_3$. The cross term does not separate, because $\overline{wT} \neq \bar{w}\,\bar{T}$. Physically,
the moist-air specific heat depends on humidity, so a wetter stream carries more energy per kelvin
and the mixed temperature is weighted by $\md_i c_{p,\mathrm{m},i}$, not by $\md_i$. Under the
code's constant-$c_p$ form, however, that weighting can be written out in closed form; using
\eqref{eq:psy:cp_moist} and $w_3$ from \eqref{eq:psy:mix_w},
\begin{equation}
T_3 \;=\; \frac{\md_1 c_{p,\mathrm{m}}(w_1)\,T_1 + \md_2 c_{p,\mathrm{m}}(w_2)\,T_2}
{(\md_1 + \md_2)\,c_{p,\mathrm{m}}(w_3)} . \label{eq:psy:mix_T_exact}
\end{equation}
So the honest statement is not that temperature is never averaged, but this: \emph{under the
constant-$c_p$ enthalpy of \eqref{eq:psy:h_coded} the mixed temperature is exactly the $\md_i
c_{p,\mathrm{m},i}$-weighted average of the inlet temperatures}, degenerating to the plain mass
average only when $w_1 = w_2$. Mixing is linear in $(w,h)$ and non-linear in $(w,T)$. Had either
specific heat been temperature-dependent, even \eqref{eq:psy:mix_T_exact} would not close in one
step and the junction would need iteration; the constant-$c_p$ assumption buys an exact,
non-iterative mixing node. The error from naive mass-averaging scales as $1.86\,(w_1-w_2)$ against
$1.006$, and at the scavenging mixer, where exhaust near \SI{19.9}{\gram\per\kilogram} meets outdoor
air near \SI{21.4}{\gram\per\kilogram}, it is hundredths of a kelvin --- a property of this
operating point, not a licence. On a dry day it grows quickly.

\paragraph{Fractional form as it appears in the code.}
At both mixing nodes the streams are fractions of one loop flow summing to unity. Writing
$d_{\mathrm{A}}$ for the fresh-air fraction and $d_{\mathrm{B}} = 1 - d_{\mathrm{A}}$ for the
recirculated one, \eqref{eq:psy:mix_w} and \eqref{eq:psy:mix_h} collapse to convex combinations with
$\md_3 = \md$:
\begin{equation}
w_3 = d_{\mathrm{B}} w_1 + d_{\mathrm{A}} w_2 , \qquad
h_3 = d_{\mathrm{B}} h_1 + d_{\mathrm{A}} h_2 , \qquad T_3 = T(h_3, w_3) . \label{eq:psy:mix_frac}
\end{equation}
This is literally the body of \texttt{process\_inlet\_state} (recirculated room return mixed with
outdoor makeup, weights $d_{\mathrm{pB}}$ and $d_{\mathrm{pA}}$) and of the scavenging node inside
\texttt{regen\_air\_pass} (building exhaust at the room state mixed with outdoor air, weights
$d_{\mathrm{rA}}$ and $d_{\mathrm{rB}}$). Both then add a fan temperature rise, set to
\SI{0}{\kelvin} in the shipped configuration, so both fans are currently isenthalpic and
\eqref{eq:psy:mix_frac} is the entire node.

\subsection{Sensible processes: what a coil does on the chart}
\label{sec:psy:sensible}

A coil that heats or cools air without condensation transfers energy but no mass:
\begin{equation}
w_{\mathrm{out}} \;=\; w_{\mathrm{in}} \qquad \text{(sensible process)} , \label{eq:psy:sensible}
\end{equation}
and the state moves \emph{horizontally} on a chart with $T$ on the abscissa and $w$ on the ordinate.
Two components are sensible by construction here: the regeneration coil, whose $\varepsilon$--NTU
routine returns only an air outlet temperature, and the air-to-air cooler, which does the same.

Heating at constant $w$ collapses the relative humidity: in \eqref{eq:psy:RH_from_w} the numerator
is fixed while $\psat(T)$ climbs steeply. This is the entire purpose of the regeneration coil --- it
does not dry the air, it makes the air \emph{thirsty}, and that collapsed relative humidity is what
pulls water out of the sorbent. The wheel model says so explicitly: its diagnostic warns that the
wheel dries the process stream only while the process-side relative humidity exceeds the
regeneration-side value, and reverses to \emph{humidifying} once that gap closes.

Cooling at constant $w$ does the opposite: $\RH$ rises, and \eqref{eq:psy:sensible} holds only while
the outlet stays above the stream's dew point. \emph{The cooler routine does not test this.} It
back-solves a conductance to hit a requested outlet temperature and returns $w_{\mathrm{out}} =
w_{\mathrm{in}}$ regardless, capping at a maximum $\UA$ if the target is unreachable. The only
supersaturation check anywhere is in the load interface, which evaluates \eqref{eq:psy:RH_from_w} on
the \emph{required} supply state and warns if it exceeds unity. That is a check on demand, one step
downstream, and it will not fire when the cooler is merely asked for an outlet close to saturation.
Any run with a low supply temperature should be audited by evaluating \eqref{eq:psy:RH_from_w} at
the cooler outlet.

\subsection{Wet-bulb temperature}
\label{sec:psy:wetbulb}

The notebook contains \emph{no} wet-bulb routine, and no dew-point routine either. This is worth
stating plainly, because a reader from a conventional HVAC background will look for one. Nothing in
the solution path needs it: cooling here is dry-coil sensible, dehumidification is by sorption
rather than condensation, and the state is carried as $(T,w)$ throughout.

Where a wet-bulb variable would appear classically, the wheel model uses Jurinak's characteristic
potentials instead. $F_1$ is very nearly constant along a line of constant enthalpy --- the
notebook's own self-check shows this by evaluating it at three temperatures along $h =
\SI{85}{\kilo\joule\per\kilogram}$ --- and constant enthalpy is, to good approximation, constant
thermodynamic wet bulb; $F_2$ plays the same role for relative humidity. The wheel therefore works
in an enthalpy-like and an $\RH$-like coordinate, carrying what a wet-bulb/dry-bulb pair would
carry, in the variables its effectiveness is actually fitted in. Should a wet-bulb temperature ever
be wanted for reporting, it is the root of the adiabatic-saturation equation solved with the same
$\psat$; that routine would have to be written, and no result here depends on one.

\subsection{Worked example: the shipped design point}
\label{sec:psy:example}

Table~\ref{tab:psy:design} evaluates the two anchor states of the model with the configuration
constants as shipped: $\patm = \SI{101325}{\pascal}$, ambient \SI{31}{\celsius} at
\SI{75}{\percent}~RH (Gulf design day), comfort \SI{37}{\celsius} at \SI{50}{\percent}~RH (hatchery
setter).

\begin{table}[htbp]
\centering
\caption{Anchor states from the shipped configuration, evaluated with
Eqs.~\eqref{eq:psy:psat}, \eqref{eq:psy:w_from_RH} and \eqref{eq:psy:h_coded}.}
\label{tab:psy:design}
\begin{tabular}{@{}lcccc@{}}
\toprule
State & $T$ [\si{\celsius}] & $\RH$ [--]
& $w$ [\si{\gram\per\kilogram}] & $h$ [\si{\kilo\joule\per\kilogram}] \\
\midrule
Outdoor air (design day)  & 31.0 & 0.75 & 21.36 & 85.8 \\
Comfort (hatchery setter) & 37.0 & 0.50 & 19.86 & 88.3 \\
\bottomrule
\end{tabular}
\end{table}

Check the outdoor value by hand: $\psat(31) = 610.94\,\exp(17.625 \times 31/274.04) =
\SI{4486}{\pascal}$; $p_{\mathrm{v}} = 0.75 \times 4486 = \SI{3365}{\pascal}$; $w = 0.622 \times
3365/(101325 - 3365) = \SI{0.02136}{\kilogram\per\kilogram}$; $h = 1.006(31) + 0.02136\,(2501 + 1.86
\times 31) = 31.2 + 54.7 = \SI{85.8}{\kilo\joule\per\kilogram}$.

Now read the table the way the model does. The outdoor air is \SI{6}{\kelvin} \emph{cooler} than the
room and its relative humidity is 25 percentage points higher, yet it carries only
\SI{1.5}{\gram\per\kilogram} more water: ventilation imposes a real but modest dehumidification duty
while contributing sensible \emph{cooling}, an unusual combination that shapes the whole control
problem, since the dominant sensible load comes from the eggs and the envelope. Note also that the
two enthalpies differ by only \SI{2.5}{\kilo\joule\per\kilogram} across that \SI{6}{\kelvin} gap ---
almost all of the sensible difference is repaid by the latent one. Reading enthalpy alone would
suggest these are nearly the same air. They are not, and that is exactly why $w$ is carried
separately.
